\subsection{And-Inverter Graphs}
\label{sec:prelim:aig}

\subsubsection{Definition and Structure}
\label{sec:prelim:aig:definition}
An And-Inverter Graph (AIG) is a Directed Acyclic Graph (DAG) $G = (\mathcal{V}, \mathcal{E})$
representing a Boolean network over the basis $\{\wedge, \neg\}$
\citep{mishchenko2006rewriting}. Every internal node is a two-input AND gate, and every edge
carries an optional inversion, so conjunction is the only operation that costs a node and
negation costs none. The basis is functionally complete, which is what allows an arbitrary
Boolean network to be expressed in one uniform structure rather than in a gate library.

The vocabulary used throughout follows from the edge set. The fanin and fanout of a node are
\begin{equation}
    \operatorname{fi}(v) = \{\, u \in \mathcal{V} : (u, v) \in \mathcal{E} \,\},
    \qquad
    \operatorname{fo}(v) = \{\, w \in \mathcal{V} : (v, w) \in \mathcal{E} \,\},
    \label{eq:prelim:fanin}
\end{equation}
with $|\operatorname{fi}(v)| = 2$ at every AND gate, $\operatorname{fi}(v) = \emptyset$ at
every Primary Input (PI) and at the constant, and $|\operatorname{fi}(v)| = 1$ at every
Primary Output (PO). Acyclicity yields a topological order $\prec$ on $\mathcal{V}$
with $u \prec v$ for every $(u, v) \in \mathcal{E}$, so signals flow from the PIs to the
POs and never return.

The POs are materialised as nodes here, rather than as pointers into the AND network as is
also common. That choice is why the node role defined below takes four values instead of
three, and it puts the circuit interface inside $|\mathcal{V}|$, the quantity the target
of~\ref{sec:method:data:label:definition} is measured on.

\subsubsection{Node Types and Inverted Edges}
\label{sec:prelim:aig:nodes}
Two labellings are carried on every graph. A node role
\begin{equation}
    \tau \colon \mathcal{V} \rightarrow
    \{\, \mathrm{const},\ \mathrm{PI},\ \mathrm{AND},\ \mathrm{PO} \,\}
    \label{eq:prelim:role}
\end{equation}
separates the constant, the PIs, the internal AND gates and the POs,
and an inversion flag $\mathrm{inv} \colon \mathcal{E} \rightarrow \{0, 1\}$ marks each edge
as a regular connection or an inverted one. Inversion is semantic rather than decorative.
Writing $b_u \in \{0,1\}$ for the value produced at $u$, the value an edge $e = (u, v)$
delivers to $v$ is
\begin{equation}
    b_u \oplus \mathrm{inv}(e),
    \label{eq:prelim:inversion}
\end{equation}
so two edges agreeing in their endpoints and differing in $\mathrm{inv}$ carry complementary
signals. The one-hot encodings of $\tau$ and $\mathrm{inv}$ are the four-valued node feature
and two-valued edge feature of~\ref{sec:method:data:representation:features}, and are not
restated there.

\subsubsection{Structural Properties: Levels, Logic Cones, Fanout}
\label{sec:prelim:aig:properties}
The topological level of a node is its longest distance from the input frontier,
\begin{equation}
    \ell(v) =
    \begin{cases}
        0 & \text{if } \operatorname{fi}(v) = \emptyset, \\[2pt]
        1 + \max_{u \in \operatorname{fi}(v)} \ell(u) & \text{otherwise,}
    \end{cases}
    \label{eq:prelim:level}
\end{equation}
and the depth of the graph is $\ell(G) = \max_{v \in \mathcal{V}} \ell(v)$. Level is an
absolute coordinate that an AIG has and a general graph does not, and four constructions in
this thesis read it: the positional encoding of~\ref{sec:method:data:representation:pe}, the
level slicing and span-weighted partitioners of~\ref{sec:method:reduction:partitioning}, and
the level bands of the cone candidate in~\ref{sec:method:summary:domain}.

Writing $u \rightsquigarrow v$ for the existence of a directed path, of length zero when
$u = v$, the fanin cone (equivalently the logic or causal cone) and the fanout cone of $v$
are
\begin{equation}
    \mathcal{C}(v) = \{\, u \in \mathcal{V} : u \rightsquigarrow v \,\},
    \qquad
    \mathcal{F}(v) = \{\, w \in \mathcal{V} : v \rightsquigarrow w \,\}.
    \label{eq:prelim:cone}
\end{equation}
Both cones contain $v$ itself. The fanin cone is the set of gates on which $v$ functionally
depends, and~\eqref{eq:prelim:level} measures the longest path through it. The claim of~\ref{sec:intro:motivation:reduction}, that
generic reduction severs causal structure, is a claim about $\mathcal{C}$: deleting one edge
on a long path removes gates from the cones of everything downstream of it, whether or not the
result resembles the original under an undirected measure. The depth-bounded form
$\mathcal{C}_k$ of~\eqref{eq:method:cone} is what makes the effect measurable.

Fanout degree $|\operatorname{fo}(v)|$ separates gates whose output is consumed once from
gates whose output is shared. Chains of single-fanout gates carry no sharing information;
multi-fanout gates are the points at which logic is reused, and they are what constrains
rewriting. Reconvergence is the same property seen from below: $v$ is a reconvergence point
of $u$ when two internally disjoint directed paths run from $u$ to $v$. AIGs are dense in
reconvergence, because structural hashing (\ref{sec:prelim:algorithms:primitives}) shares
every duplicated subexpression instead of replicating it \citep{kuehlmann2002robust}.

Two dominance relations are needed, and only one of them is informative on an AIG\@. Add a
virtual source $s$ preceding every PI and a virtual sink $t$ succeeding every PO\@. A node
$u$ dominates $v$ when every path $s \rightsquigarrow v$ passes through $u$, and
post-dominates $v$ when every path $v \rightsquigarrow t$ passes through $u$; the immediate
dominator $\operatorname{idom}(v)$ and immediate post-dominator
$\operatorname{ipdom}(v)$ are the closest such nodes \citep{lengauer1979dominators}. A gate below the input frontier is
reachable from $s$ along many independent paths, so $\operatorname{idom}(v)$ degenerates to
$s$ for most AND gates and partitions nothing. The post-dominator does not degenerate the
same way. Here $\operatorname{ipdom}(v)$ is the gate at which the fanout cone
$\mathcal{F}(v)$ reconverges, or $t$ when that cone never reconverges, and grouping by it
is what the cone candidate of~\ref{sec:method:summary:domain} uses.
% TODO(measure): idom degeneracy asserted, not yet measured on this corpus.

The Maximal Fanout-Free Cone (MFFC) of a root gate $v$ collects the gates feeding $v$ and
nothing else \citep{mishchenko2006rewriting},
\begin{equation}
    \operatorname{MFFC}(v) = \{\, u \in \mathcal{C}(v) :
        \text{every path } u \rightsquigarrow t \text{ passes through } v \,\},
    \label{eq:prelim:mffc}
\end{equation}
so removing $v$ from the network removes exactly $\operatorname{MFFC}(v)$ and nothing more.
That property makes the MFFC the natural unit of local restructuring, and refactoring
collapses and re-expresses one MFFC at a time
(\ref{sec:prelim:algorithms:primitives}). It is also the unit merged by the second domain
candidate of~\ref{sec:method:summary:domain}, which is why that candidate carries an argument
about label leakage and the others do not.

\subsubsection{Why AIGs Differ from General Graphs}
\label{sec:prelim:aig:distinct}
Four properties separate an AIG from the citation, social and knowledge graphs on which graph
reduction was developed, and each costs a generic method something specific.

(i) An AIG is directed and acyclic. Machinery built for undirected graphs does not transfer
unchanged, because a quotient of a DAG need not be acyclic, and a reduction computed on the
undirected projection discards the orientation that carries the logic. The spanning forest
of~\ref{sec:method:sparse:forest} is the concrete case.

(ii) An AIG is deep relative to any practical encoder. Against an encoder of $L$
message-passing layers, a node representation sees an $L$-hop prefix of a fanin cone whose
length is $35$ levels at the corpus median and $24{,}802$ at its maximum
(\ref{sec:method:data:sources}). Receptive-field starvation is therefore the normal condition
on this data, and it is the starvation that hypothesis H1 (\ref{sec:intro:rq3}) proposes
coarsening relieves.

(iii) Edges carry function. The inversion in~\eqref{eq:prelim:inversion} is part of what the
graph represents, so a merge that averages or discards $\mathrm{inv}$ changes the Boolean function rather
than approximating it. Every summarization method here is therefore constrained to keep
normal and inverted connections distinguishable, a constraint no general method imposes on
itself.

(iv) The interface is fixed. The PIs and POs are the circuit's contract with the
rest of the design, and merging them away changes what the graph is rather than coarsening
it. The boundary is enforced uniformly in~\ref{sec:method:summary:mergemap}.

A reduction operator can respect or ignore all four at no syntactic cost, since the graph it
produces is well formed either way. Whether respecting them buys accuracy is the question RQ4
(\ref{sec:intro:rq4}) measures.
